\section{Planar equivalent-layer theorems}
\label{sec:solutions_eqlayer_integral}

\subsection{Preliminaries}

Let $d^{o}(\mathbf{r})$ be the harmonic function defining a given 
observed potential field at the observation point 
$\mathbf{r} \notin \overline{\mathcal{V}}$ outside the sources. 
Let us also consider that this
potential field is defined in terms of the potential $U(\mathbf{r})$
(equation \ref{eq:scalar-potential-U}) so that its 2DFT on a horizontal plane
$z^{\dagger}_{0} < z^{t}_{0}$ above the sources is given by:
\begin{equation}
	\mathcal{F} [ d^{o}(\mathbf{r}^{\dagger}) ] =
	\alpha_{U} \mathcal{F} [ U(\mathbf{r}^{\dagger}) ] \: ,
	\label{eq:Fourier-observed-data}
\end{equation}
where $\alpha_{U}$ defines the relationship between $d^{o}(\mathbf{r})$ 
and $U(\mathbf{r})$ as a function of the wavenumbers $k_{x}$, $k_{y}$ and 
$k = \sqrt{k_{x}^{2} + k_{y}^{2}}$. 
Here, we consider two kinds of $\alpha_{U}$ given by
\begin{equation}
	\alpha_{U} = (\mathtt{i} \, k_{x})^{\ell} (\mathtt{i} \, k_{y})^{m} k^{n} \: ,
	\quad \ell, m , n \in \{0, 1, 2\} \: ,
	\label{eq:alpha_U_monopoles}
\end{equation}
or
\begin{equation}
	\alpha_{U} = (\mathbf{k} \cdot \hat{\mathbf{h}})^{\ell} \, (\mathbf{k} \cdot \hat{\mathbf{h}})^{m} \: ,
	\quad \ell, m \in \{0, 1\} \: ,
	\label{eq:alpha_U_dipoles}
\end{equation}
where $\hat{\mathbf{h}}$ and $\hat{\mathbf{t}}$ are unit vectors and 
$\mathbf{k} = ( \mathtt{i} \, k_{x} \, , \, \mathtt{i} \, k_{y} \, , \, k)^{\top}$,
according to the compact notation proposed by \citet{pedersen1978}.
For the subsequent developments, it is important to draw attention to the fact that 
$\alpha_{U}$ (equations \ref{eq:alpha_U_monopoles} and \ref{eq:alpha_U_dipoles}) admit 
the upper bound
\begin{equation}
	\alpha_{U} \le (3 \, k)^{6} \: .
	\label{eq:alpha_U_upper_bound}
\end{equation}

Using the upper bound for the amplitude spectra of $U(\mathbf{r}^{\dagger})$
(equation \ref{eq:abs-Fourier-U-r-dagger-bound}) and the 2DFT of $d^{o}(\mathbf{r}^{\dagger})$
(equation \ref{eq:Fourier-observed-data}), we deduce
\begin{equation}
	\vert \mathcal{F} [ d^{o}(\mathbf{r}^{\dagger}) ] \vert \le
	\vert \alpha_{U} \vert \, \Omega \, 
	\frac{e^{- k \, \Delta z}}{k} \: ,
	\: \quad \Delta z = z^{t}_{0} - z^{\dagger}_{0} > 0 \: .
	\label{eq:abs-Fourier-general-derivative-U-bound}
\end{equation}
Note that, by using $\alpha_{U} = k$, we obtain the same result
presented by \citet{bhattacharyya1967} as the upper bound for the amplitude spectrum of 
$\partial_{z} U(\mathbf{r})$ (equation \ref{eq:1st-deriv-U}), whereas 
$\alpha_{U} = 1$ leads to the upper bound defined by equation \ref{eq:abs-Fourier-U-r-dagger-bound}.

Let us now apply the Fourier-convolution theorem to the planar equivalent-layer integral
(equation \ref{eq:integral-equation-eqlayer}) on a horizontal plane 
$z^{\dagger}_{0} < z^{t}_{0}$ above the sources to obtain:
\begin{equation}
	\mathcal{F} [ d(\mathbf{r}^{\dagger}, z''_{0}) ] =
	\mathcal{F} [ p(\mathbf{r}'') ] \, 
	\mathcal{F} [ a(\mathbf{r}^{\dagger}, \mathbf{r}'') ] \: ,
	\label{eq:Fourier-integral-equation-eqlayer}
\end{equation}
where
\begin{equation}
	\mathcal{F} [ a(\mathbf{r}^{\dagger}, \mathbf{r}'') ] =
	\alpha_{a} \, 
	\mathcal{F} [ \Phi(\mathbf{r}^{\dagger}, \mathbf{r}'') ] \: ,
	\label{eq:Fourier-kernel}
\end{equation}
and $\alpha_{a}$ defines the relationship between 
$a(\mathbf{r}^{\dagger}, \mathbf{r}'')$ (equation \ref{eq:integral-equation-eqlayer}) 
and $\Psi(\mathbf{r}^{\dagger}, \mathbf{r}'')$ (equation \ref{eq:Psi})
in terms of the wavenumbers $k_{x}$, $k_{y}$ and 
$k = \sqrt{k_{x}^{2} + k_{y}^{2}}$. 
Here, we limit ourselves to the particular case in which $\alpha_{a}$ is
defined in the same way as $\alpha_{U}$ (equations \ref{eq:alpha_U_monopoles} 
and \ref{eq:alpha_U_dipoles}).
%Our next step consists in 
%defining an equivalent layer with physical property distribution 
%$\mathcal{F} [ p(\mathbf{r}'') ]$ that produces a
%predicted potential field $\mathcal{F} [ d(\mathbf{r}^{\dagger}, z''_{0}) ]$
%(equation \ref{eq:Fourier-integral-equation-eqlayer}) close to the observed
%potential field $\mathcal{F} [ d^{o}(\mathbf{r}^{\dagger}) ]$
%(equation \ref{eq:Fourier-observed-data}).


%\subsection{Exact solutions on a horizontal plane above the sources}
\subsection{Planar equivalent-layer theorem on a horizontal plane above the sources}
\label{subsec:exact_eqlayer_integral}

Consider a horizontal plane $z''_{0} < z^{t}_{0}$ located above the sources.
In this case, there exist solutions $p(\mathbf{r}'')$ 
of the planar equivalent-layer integral (equation \ref{eq:integral-equation-eqlayer})
producing predicted potential fields $d(\mathbf{r}, z''_{0})$
that exactly reproduce the observed potential field $d^{o}(\mathbf{r}^{\dagger})$
at any point $\mathbf{r}$ with $z < z''_{0}$.

A proof of this theorem can be obtained by first defining 
the 2DFT of $p(\mathbf{r}'')$ on the horizontal plane $z''_{0}$ in terms 
of the exact solutions of the Dirichlet's and Neumann's problems 
in Fourier-domain (equation \ref{eq:Fourier-U-Dirichlet-Neumann}) as follows:
\begin{equation}
	\mathcal{F} [ p(\mathbf{r}'') ] =
	\tfrac{k \, \alpha_{U}}{\alpha_{a} \, \alpha_{\Phi}} \,
	\mathcal{F} [ U(\mathbf{r}'') ] 
	\: , \quad z''_{0} < z^{t}_{0} \: .
	\label{eq:Fourier-p-exact}
\end{equation}
Then, by substituting $\mathcal{F} [ a(\mathbf{r}^{\dagger}, \mathbf{r}'') ]$ 
(equation \ref{eq:Fourier-kernel}) and $\mathcal{F} [ p(\mathbf{r}'') ]$
(equation \ref{eq:Fourier-p-exact}) into the planar equivalent-layer integral
in Fourier domain (equation \ref{eq:Fourier-integral-equation-eqlayer}) we obtain
\begin{equation}
	\mathcal{F} [ d(\mathbf{r}^{\dagger}, z''_{0}) ] = 
	\alpha_{U} \, \mathcal{F} [ U(\mathbf{r}^{\dagger}) ] 
	\: , \quad z^{\dagger}_{0} < z''_{0} < z^{t}_{0} \: .
	\label{eq:Fourier-predicted-data-exact}
\end{equation}
From equations \ref{eq:Fourier-observed-data} and \ref{eq:Fourier-predicted-data-exact},
we finally have that 
\begin{equation}
	\mathcal{F} [ d^{o}(\mathbf{r}^{\dagger}) ] =
	\mathcal{F} [ d(\mathbf{r}^{\dagger}, z''_{0}) ]
	\: , \quad z^{\dagger}_{0} < z''_{0} < z^{t}_{0} \: ,
	\label{eq:Fourier-datafit-exact}
\end{equation}
and, consequently, that 
$d^{o}(\mathbf{r}) = d(\mathbf{r}, z''_{0})$ at any point 
$\mathbf{r}$ with $z < z''_{0} < z^{t}_{0}$.


%\subsection{Approximate solutions on a horizontal plane crossing the sources}
\subsection{Planar equivalent-layer theorem on a horizontal plane crossing the sources}
\label{subsec:approx_eqlayer_integral}

Consider a horizontal plane $z''_{0} > z^{t}_{0}$ that crosses the sources.
In this case, there exist solutions $p(\mathbf{r}'')$ 
of the planar equivalent-layer integrals (equation \ref{eq:integral-equation-eqlayer})
producing predicted potential fields $d(\mathbf{r}, z''_{0})$
that approximate the observed potential field $d^{o}(\mathbf{r})$ to 
a predefined error $\eta > 0$ at any horizontal plane $z^{\dagger}_{0} < z^{t}_{0}$
so that
\begin{equation}
	\int\limits_{-\infty}^{\infty} \int\limits_{-\infty}^{\infty}
	\vert d^{o}(\mathbf{r}^{\dagger}) - 
	d(\mathbf{r}^{\dagger}, z''_{0}) \vert^{2}
	\:\: dx^{\dagger}dy^{\dagger} \le \eta 
	\: , \quad z^{\dagger}_{0} < z^{t}_{0} < z''_{0} \: ,
	\label{eq:eta-bound-data}
\end{equation}
where $\eta > 0$ is a fixed but arbitrarily small error level.

This theorem can be proved by using the general approximate solution of the Dirichlet's and
Neumann's problems (equation \ref{eq:Fourier-tau}) to define a
$\mathcal{F} [ p(\mathbf{r}'') ]$ given by
\begin{equation}
	\mathcal{F} [ p(\mathbf{r}'') ] =
	\tfrac{k \, \alpha_{U}}{\alpha_{a} \, \alpha_{\Phi}} \,
	\mathcal{F} [ \tau(\mathbf{r}'') ] 
	\: , \quad z^{t}_{0} < z''_{0} \: ,
	\label{eq:Fourier-p-approx}
\end{equation}
which combined with equations \ref{eq:Fourier-integral-equation-eqlayer} and
\ref{eq:Fourier-kernel} results in
\begin{equation}
	\mathcal{F} [ d(\mathbf{r}^{\dagger}, z''_{0}) ] = 
	\alpha_{U} \, \mathcal{F} [ V(\mathbf{r}^{\dagger}, z''_{0}) ] 
	\: , \quad z^{\dagger}_{0} < z^{t}_{0} < z''_{0} \: .
	\label{eq:Fourier-predicted-data-approx}
\end{equation}
Next we apply the Plancherel's theorem \citep[][ p. 119]{bracewell1999} to the 
left-hand side of equation \ref{eq:eta-bound-data}, use 
$\mathcal{F} [ V(\mathbf{r}^{\dagger}, z''_{0})$ (equation \ref{eq:Fourier-V-k-epsilon}),
$\mathcal{F} [ d^{o}(\mathbf{r}^{\dagger}) ]$ (equation \ref{eq:Fourier-observed-data})
and $\mathcal{F} [ d(\mathbf{r}^{\dagger}, z''_{0}) ]$ (equation \ref{eq:Fourier-predicted-data-approx})
and change variables with $k_{x} = k \, \cos\theta$, $k_{y} = k \, \sin\theta$ and 
$k = \sqrt{k_{x}^{2} + k_{y}^{2}}$ to obtain
\begin{equation}
	\int\limits_{-\infty}^{\infty} \int\limits_{-\infty}^{\infty}
	\vert 
	d^{o}(\mathbf{r}^{\dagger}) - d(\mathbf{r}^{\dagger}, z''_{0})
	\vert^{2} 
	\: dx^{\dagger}dy^{\dagger} =
	\int\limits_{0}^{2\pi} \int\limits_{k_{\varepsilon}}^{\infty}
	\vert \mathcal{F} [ d^{o}(\mathbf{r}^{\dagger}) ] \vert^{2} 
	\: k \: d k \, d \theta \: .
	\label{eq:Plancherel-k-theta-data}
\end{equation}
Then we define a conservative $\eta$ (equation \ref{eq:eta-bound-data})
by substituting the upper bounds for $\alpha_{U}$ (equation \ref{eq:alpha_U_upper_bound}) 
and for the amplitude spectrum of $d^{o}(\mathbf{r}^{\dagger})$ 
(equation \ref{eq:abs-Fourier-general-derivative-U-bound}) into the right-hand side of 
equation \ref{eq:Plancherel-k-theta-data} as follows
\begin{equation}
	\eta \equiv 2\pi \, \Omega^{2} \,
	\int\limits_{k_{\varepsilon}}^{\infty}
	\frac{e^{- 2 \, \Delta z \, k }}{k} \, (3 \, k)^{6}
	\: d k \: .
	\label{eq:eta}
\end{equation}
Finally, we use equation \ref{eq:amplitude-spectrum-U-inequality} to show that
\begin{equation}
	\int\limits_{0}^{2\pi} \int\limits_{k_{\varepsilon}}^{\infty}
	\vert \mathcal{F} [ d^{o}(\mathbf{r}^{\ddagger}) ] \vert^{2} \: k \, dkd\theta <
	\int\limits_{0}^{2\pi} \int\limits_{k_{\varepsilon}}^{\infty}
	\vert \mathcal{F} [ d^{o}(\mathbf{r}^{\dagger}) ] \vert^{2} \: k \, dkd\theta \: ,
	\quad z^{\ddagger}_{0} < z^{\dagger}_{0} \: ,
	\label{eq:amplitude-spectrum-do-inequality}
\end{equation}
and, consequently, that 
$d(\mathbf{r}, z''_{0}) \approx d^{o}(\mathbf{r})$ to the 
predefined error $\eta$ (equation \ref{eq:eta}) according to
equation \ref{eq:eta-bound-data} not only on the 
horizontal plane $z^{\dagger}_{0}$, but also at any horizontal plane above it,
with $z < z^{\dagger}_{0} < z^{t}_{0} < z''_{0}$.

\subsection{Remarks on particular solutions of the planar equivalent-layer integral}
\label{subsec:remarks_eqlayer_integral}

The Fourier-domain solutions $\mathcal{F} [ p(\mathbf{r}'') ]$ associated with equivalent layers located
above the sources (equation \ref{eq:Fourier-p-exact}) and that cross the sources
(equation \ref{eq:Fourier-p-approx}) are essentially the same, depending on the truncating
wavenumber $k_{\varepsilon}$ (equation \ref{eq:Fourier-tau}). 
The physical meaning of a given equivalent layer is defined by the respective $\alpha_{\Phi}$
(equation \ref{eq:alpha-Phi}), $\alpha_{U}$ (equation \ref{eq:Fourier-observed-data}) 
and $\alpha_{a}$ (equation \ref{eq:Fourier-kernel}).

\subsubsection{Solutions defined by $\alpha_{\Phi} = k$ and $\alpha_{a} = 1$}

All equivalent layers obtained with $\alpha_{\Phi} = k$ and $\alpha_{a} = 1$ 
give rise to a predicted potential field $d(\mathbf{r}, z''_{0})$ which is analogous
to the vertical component of the gravitational attraction due to a layer of monopoles.
In this case, the physical property distribution $p(\mathbf{r}'')$ is the observed 
potential field $d^{o}(\mathbf{r})$ (or its filtered version) continued to the layer plane
and, consequently, the planar equivalent-layer integral (equation \ref{eq:integral-equation-eqlayer})
represents the well-known \textit{upward continuation integral} 
\citep[e.g.,][]{skeels1947, peters1949, henderson-zietz1949, dean1958, roy1962, bhattacharyya1967, henderson1970},
which is the exact or approximate solution of the Dirichlet’s problem on the
layer plane \cite[][ p. 236]{kellogg1967}.

In the Fourier domain, the equivalent-layer integral assumes the form:
%\begin{equation}
%	
%	\label{eq:eqlayer-integral-dirichlet-Fourier}
%\end{equation}

 and has a kernel 
$a(\mathbf{r}^{\dagger}, \mathbf{r}'')$ given by 
$\mathcal{F} [ a(\mathbf{r}^{\dagger}, \mathbf{r}'') ] = k \, \mathcal{F} [ \Psi(\mathbf{r}^{\dagger}, \mathbf{r}'') ]$
(equations \ref{eq:Fourier-Phi}--\ref{eq:Fourier-Psi}) in the Fourier domain or
$a(\mathbf{r}^{\dagger}, \mathbf{r}'') = \partial_{z} \Psi(\mathbf{r}^{\dagger}, \mathbf{r}'')$
(equation \ref{eq:Psi}) in the space domain.

Note that such equivalent layers preserve the physical relationship between the 
observed potential field and its true sources only for the particular case in which 
the observed potential field is gravity disturbance data \cite[e.g.,][ p. 93]{hofmann-wellenhof-moritz2005}, 
i.e., for the case in which $d^{o}(\mathbf{r}) = \partial_{z} U(\mathbf{r})$ (equation \ref{eq:1st-deriv-U}).
The solution of the planar equivalent-layer integral (equation \ref{eq:integral-equation-eqlayer})
obtained with $\alpha_{\Phi} = k$ and $\alpha_{a} = 1$ is the basis of the pioneer method
proposed by \cite{dampney1969} for processing gravity data.

\subsubsection{Solutions defined by $\alpha_{\Phi} = 1$ and $\alpha_{a} = \alpha_{U}$}

Most of what is known or assumed to be true without any proof about planar equivalent 
layers can be deduced from the set of solutions $p(\mathbf{r}'')$ of the planar equivalent-layer 
integral (equation \ref{eq:integral-equation-eqlayer}) obtained with 
$\alpha_{\Phi} = 1$ and $\alpha_{a} = \alpha_{U}$,
which are directly related to the solution of the Neumann’s problem for the potential $U(\mathbf{r})$ 
(equation \ref{eq:scalar-potential-U}) on the layer plane \cite[][ p. 246]{kellogg1967}. 
All these equivalent layers preserve the physical relationship between the 
observed potential field and its true sources.
Besides, all these equivalent layers have physical property distribution $p(\mathbf{r}'')$ defined 
by the vertical derivative of $U(\mathbf{r})$ (equation \ref{eq:scalar-potential-U})
or $V(\mathbf{r}, z''_{0})$ (equation \ref{eq:potential-V}) on the layer plane,
regardless of how $\alpha_{U}$ is set.

The set of solutions $p(\mathbf{r}'')$ of the planar equivalent-layer 
integral (equation \ref{eq:integral-equation-eqlayer}) obtained with 
$\alpha_{\Phi} = 1$ and $\alpha_{a} = \alpha_{U}$ results in two groups of equivalent 
layers, according to $\alpha_{U}$.
The first group has $\alpha_{U}$ given by equation \ref{eq:alpha_U_monopoles} and is formed 
by layers of monopoles.
In this case, the observed potential field $d^{o}(\mathbf{r}^{\dagger})$
(equation \ref{eq:Fourier-observed-data}) represents the potential $U(\mathbf{r})$ 
(equation \ref{eq:scalar-potential-U}) or its derivatives 
(equations \ref{eq:1st-deriv-U} and \ref{eq:2nd-deriv-U}), which describe the gravitational 
field produced by the true sources.
It can be shown that any given layer of monopoles in this group, defined 
on a specific horizontal plane $z''_{0}$, simultaneously reproduces (or approximates) the
potential $U(\mathbf{r})$ (equation \ref{eq:scalar-potential-U}) and all its derivatives
(equations \ref{eq:1st-deriv-U} and \ref{eq:2nd-deriv-U}).
This result gives the theoretical support for jointly processing gravity-gradient data 
\cite[e.g.,][]{barnes_lumley2011} or converting gravity disturbance data into the 
gravity-gradient tensor \cite[e.g.,][]{piauilino_etal2019_seg} via equivalent-layer 
technique, for example.
We draw attention to the fact that these potential-field applications implicitly presume
the existence of a single layer of monopoles that satisfactorily reproduces all components
of the gravitational field. This premise, however, is usually assumed
to be true without any proof.

The second group of solutions $p(\mathbf{r}'')$ obtained with 
$\alpha_{\Phi} = 1$ and $\alpha_{a} = \alpha_{U}$ is formed by layers of dipoles.
In this case, the
observed potential field $d^{o}(\mathbf{r}^{\dagger})$ (equation \ref{eq:Fourier-observed-data}) 
represents directional derivatives of the potential $U(\mathbf{r})$ 
(equations \ref{eq:oblique-1st-deriv-U} and \ref{eq:oblique-2nd-deriv-U}),
which describe the magnetic induction field due to the true sources.

A PARTE ABAIXO SE REFERE A FONTES COM DIREÇÃO UNIFORME E CAMADA COM A MESMA
DIREÇÃO DA FONTE - ISSO PRECISA SER ESPECIFICADO

It can be shown that any given layer of dipoles in this group, defined 
on a specific horizontal plane $z''_{0}$, simultaneously reproduces 
(or approximates) all the directional derivatives of the
potential $U(\mathbf{r})$ (equations \ref{eq:oblique-1st-deriv-U} and 
\ref{eq:oblique-2nd-deriv-U}).

This result is important because most previously published studies are limited to
the empirical use of a planar layer of dipoles for processing total-field anomaly data.

Few studies \cite[][]{pedersen1991, li_etal2014, reis_etal2020} have drawn
attention to the problem of proving the existence of a planar layer of dipoles
that reproduces the approximated total-field anomaly $\partial_{th} U(\mathbf{r})$
(equation \ref{eq:oblique-2nd-deriv-U}).



FALAR DA IMPORTANCIA DESTE RESULTADO PARA O PROCESSAMENTO DE DADOS MAGNETICOS, 
ESTIMATIVA DAS COMPONENTS DO CAMPO VIA AN DE CAMPO TOTAL APROX E RTP

%For example, defining $\alpha_{U}$ (equation \ref{eq:alpha_U_monopoles}) with $\ell = m = 0$ 
%and $n = 1$ leads to an observed potential field 
%$\mathcal{F} [ d^{o}(\mathbf{r}^{\dagger}) ] = k \, \mathcal{F} [ U(\mathbf{r}^{\dagger}) ]$
%(equation \ref{eq:Fourier-observed-data}) and a kernel
%$\mathcal{F} [ a(\mathbf{r}^{\dagger}, \mathbf{r}'') ] = k \, \mathcal{F} [ \Psi(\mathbf{r}^{\dagger}, \mathbf{r}'') ]$
%(equations \ref{eq:Fourier-Phi}--\ref{eq:Fourier-Psi}) in the Fourier domain, or 
%$d^{o}(\mathbf{r}) = \partial_{z} U(\mathbf{r})$ (equation \ref{eq:1st-deriv-U}) and
%$a(\mathbf{r}^{\dagger}, \mathbf{r}'') = \partial_{z} \Psi(\mathbf{r}^{\dagger}, \mathbf{r}'')$
%(equation \ref{eq:Psi}) in the space domain.
%If we now use $\ell = m = 1$ and $n = 0$ (equation \ref{eq:alpha_U_monopoles}), we obtain
%$\mathcal{F} [ d^{o}(\mathbf{r}^{\dagger}) ] =  -k_{x} k_{y} \, \mathcal{F} [ U(\mathbf{r}^{\dagger}) ]$
%(equation \ref{eq:Fourier-observed-data}) and a kernel
%$\mathcal{F} [ a(\mathbf{r}^{\dagger}, \mathbf{r}'') ] = -k_{x} k_{y} \, \mathcal{F} [ \Psi(\mathbf{r}^{\dagger}, \mathbf{r}'') ]$
%(equations \ref{eq:Fourier-Phi}--\ref{eq:Fourier-Psi}) in the Fourier domain,
%or $d^{o}(\mathbf{r}) = \partial_{xy} U(\mathbf{r})$ (equation \ref{eq:2nd-deriv-U}) and
%$a(\mathbf{r}^{\dagger}, \mathbf{r}'') = \partial_{xy} \Psi(\mathbf{r}^{\dagger}, \mathbf{r}'')$
%(equation \ref{eq:Psi}) in the space domain.